\section{A small obstruction to eliminating the offset parameter}
The dominated implicational axiom is
\[
 \rho(D,A,H,u)=((H\to u)\to D)\to((D\to H)\to(A\to H)),
 \qquad u\le D.
\]
Two tempting simplifications are to keep only \(u=D\), or only \(u=0\).
The following matrices show that neither restricted family, considered
as a propositional axiom schema with modus ponens and substitution,
even proves the weakening axiom \(P\to(Q\to P)\).
This is a negative result about those proposed replacements of the
propositional interface alone; it does not show that the full numerical
calculus with all of its other axioms is incomplete after either change.

For the reflexive restriction \(u=D\), use values \(\{0,1,2\}\),
designate only \(2\), and define implication by the following table:
\[
\begin{array}{c|ccc}
\to&0&1&2\\\hline
0&1&1&2\\
1&1&2&2\\
2&1&1&2
\end{array}.
\]
For the restriction \(u=0\), interpret the source truth constant
zero by the matrix value \(2\), and use
\[
\begin{array}{c|ccc}
\to&0&1&2\\\hline
0&0&0&2\\
1&2&2&2\\
2&1&1&2
\end{array}.
\]
Both tables satisfy \(2\to2=2\), so the numerical truth alias
\(I(0,0)=0\) is respected. Modus ponens preserves designation because
the last row has its sole designated entry in column \(2\).
All 27 instances of the relevant three-variable restricted schema
evaluate to \(2\). In contrast, with \(P=Q=0\) the weakening instance
evaluates to \(1\) in the first matrix and to \(0\) in the second.
The companion script checks every displayed claim directly.

Consequently an offset-eliminating redesign cannot justify its
propositional completeness merely by restricting the old four-parameter
head and reusing its proof. It would have to derive the missing
implicational principles using other retained axiom families, or supply
a different complete basis and implementation.
